% Part: first-order-logic
% Chapter: tableaux
% Section: derivations

\documentclass[../../../include/open-logic-section]{subfiles}

\begin{document}

\iftag{FOL}
      {\olfileid{fol}{tab}{der}}
      {\olfileid{pl}{tab}{der}}

\olsection{\usetoken{P}{tableau}}

\begin{explain}
Kita wis nerangake apa iku asumsi lan wis menehi aturan inferensi.
!!^{tableau}s diasilake kanthi induktif saka unsur-unsur iki: saben
!!{tableau} iku salah siji saka cabang tunggal kang dumadi saka siji
utawa luwih asumsi, utawa diasilake saka !!a{tableau} kanthi
ngetrapake salah siji aturan inferensi ing sawijining cabang.
\end{explain}

\begin{defn}[!!^{tableau}]
!!^a{tableau} kanggo asumsi $\sFmla{S_1}{!A_1}$, \dots,
$\sFmla{S_n}{!A_n}$ (kanthi saben $S_i$ iku salah siji saka $\True$
utawa~$\False$) yaiku wit winates saka !!{signed formula}s kang
nyukupi syarat-syarat iki:
\begin{enumerate}
\item $n$ !!{signed formula}s paling dhuwur ing wit yaiku
  $\sFmla{S_i}{!A_i}$, siji ana ing sangisore liyane.
\item Saben !!{signed formula} ing wit kang dudu salah siji saka
  asumsi diasilake saka panrapan aturan inferensi kang bener marang
  !!a{signed formula} ing cabang ing sadhuwure.
\end{enumerate}
Sawijining cabang saka !!a{tableau} iku \emph{katutup} yen lan mung
yen ngemot $\sFmla{\True}{!A}$ lan~$\sFmla{\False}{!A}$ bebarengan,
lan \emph{bukak} yen ora mangkono. !!^a{tableau} kang saben cabange
katutup diarani \emph{!!{tableau} katutup} (kanggo himpunan
asumsine). Yen !!a{tableau} ora katutup, yaiku yen ngemot saora-orane
siji cabang bukak, tableau kasebut \emph{bukak}.
\end{defn}

\begin{ex}
  Saben himpunan asumsi dhewe wis dadi !!a{tableau}, nanging umume
  ora katutup. (Temenan, tableau kasebut katutup mung yen asumsi wis
  ngemot pasangan !!{signed formula}s $\sFmla{\True}{!A}$
  lan~$\sFmla{\False}{!A}$.)

  Saka !!a{tableau} (bukak utawa katutup), kita bisa ngasilake tableau
  anyar kang luwih gedhe kanthi ngetrapake salah siji aturan inferensi
  marang !!a{signed formula}~$!A$ ing jerone. Aturan kasebut bakal
  nambahake siji utawa luwih !!{signed formula}s ing pucuk saben
  cabang kang ngemot kedadeyan~$!A$ kang dikenani aturan.

  Tuladhane, gatekna asumsi $\sFmla{\True}{!A \land \lnot !A}$. Ing
  ngisor iki !!{tableau} (bukak) kang mung dumadi saka asumsi iku:
  \begin{center}
    \begin{tableau}{}
      [\sFmla{\True}{\formula{A} \land \lnot \formula{A}}, just=\TAss]
    \end{tableau}{}
  \end{center}
  Kita ngasilake !!{tableau} anyar kanthi ngetrapake aturan
  $\TRule{\True}{\land}$ marang asumsi. Aturan iku ngidini kita
  nambahake rong larik anyar ing !!{tableau},
  $\sFmla{\True}{!A}$ lan $\sFmla{\True}{\lnot !A}$:
  \begin{center}
    \begin{tableau}{}
      [\sFmla{\True}{\formula{A} \land \lnot \formula{A}}, just=\TAss
        [\sFmla{\True}{\formula{A}}, just={\TRule{\True}{\land}[1]},
          [\sFmla{\True}{\lnot \formula{A}}, just={\TRule{\True}{\land}[1]}
          ]
        ]
      ]
    \end{tableau}{}
  \end{center}
  Nalika nulis !!{tableau}s, kita nyathet aturan kang wis ditrapake
  ing sisih tengen (tuladhane, $\TRule{\True}{\land} 1$ ateges
  !!{signed formula} ing larik iku diasilake saka panrapan aturan
  $\TRule{\True}{\land}$ marang !!{signed formula} ing larik~$1$).
  !!{tableau} anyar iki saiki ngemot !!{signed formula}s tambahan,
  nanging mung siji ($\sFmla{\True}{\lnot !A}$) kang isih bisa
  dikenani aturan (ing kasus iki, aturan $\TRule{\True}{\lnot}$).
  Panrapan iku ngasilake !!{tableau} katutup iki:
  \begin{center}
    \begin{tableau}{}
      [\sFmla{\True}{\formula{A} \land \lnot \formula{A}}, just=\TAss
        [\sFmla{\True}{\formula{A}}, just={\TRule{\True}{\land}[1]}
          [\sFmla{\True}{\lnot \formula{A}}, just={\TRule{\True}{\land}[1]}
            [\sFmla{\False}{\formula{A}}, just={\TRule{\True}{\lnot}[3]}, close]
          ]
        ]
      ]
    \end{tableau}{}
  \end{center}
\end{ex}

\end{document}
